\section{成长与求学：从模拟世界到理解世界}

\subsection{童年：用加减乘除创建虚拟世界}

胡渊鸣的父母从事计算机教学工作，他在幼儿园时期就接触了计算机，玩过仙剑奇侠传、红色警戒、帝国时代等游戏。这些经历催生了一个核心感悟：\textbf{用加减乘除就能在计算机里创建一个世界}。

更深层的哲学思考在那时就已萌芽。他回忆小时候看楼下的桃树，思考一个问题："当我不在观测这棵桃树的时候，它是在默默生长，还是在偷懒保持原样，直到我看到它的一瞬间才长成应有的样子？"后来他意识到，作为被模拟的一员，你无法区分这两种情况——这与量子力学的"观测者效应"异曲同工。

\begin{knowledgebox}{模拟假说（Simulation Hypothesis）}
胡渊鸣从儿时的哲学直觉延伸到对物理世界的理解：普朗克常数可能是模拟系统的浮点数精度限制，光速可能是计算资源不足导致的交互距离限制，宇宙大爆炸可能只是初始条件设置。这种思维方式深刻影响了他后来对物理仿真和 AI 的看法。
\end{knowledgebox}

上小学后，他用 Visual Basic 6.0 编写物理仿真代码，在尚未学习牛顿定律的情况下，通过编程"自己发现"了牛顿第二定律——当他用键盘方向键直接移动物体时觉得不真实，于是给物体定义了速度增量，后来才知道这就是 $F = ma$。

\begin{table}[H]
\centering
\caption{胡渊鸣各阶段的编程与仿真经历}
\begin{tabular}{lll}
\toprule
\textbf{阶段} & \textbf{工具/语言} & \textbf{仿真内容} \\
\midrule
幼儿园/小学 & Visual Basic 6.0 & 宇宙飞船游戏、基础物理仿真 \\
小学/初中 & RPG Maker XP, Ruby & RPG游戏、弹簧质点系统(Mass-Spring) \\
高中 & C++ & 刚体仿真(Rigid Body)、多米诺骨牌 \\
本科（清华） & C++, CUDA & 流体仿真、SIGGRAPH 论文复现 \\
博士（MIT） & Taichi, LLVM & MPM、可微仿真、编译器 \\
\bottomrule
\end{tabular}
\end{table}

初中时他构建了弹簧质点系统（Mass-Spring System），灵感来自游戏 World of Goo——你有很多球，可以把它们连起来搭成建筑物，在计算机里面表现出来就是 mass-spring system。高中时开始研究刚体仿真（Rigid Body Simulation），尝试在电脑中模拟多米诺骨牌——"因为在现实世界中堆了一个东西，没有场地，而且不小心碰一下就全完蛋了。"

在这个过程中他发现：\textbf{两个刚体之间的碰撞、摩擦等模型都是近似模型}。这一认知后来深刻影响了他对 simulation 局限性的理解——也奠定了他后来在博士阶段做可微仿真、以及对 data-driven simulation 持开放态度的基础。

初中到高中期间，大量时间被计算机竞赛占用（NOI 等），competitive programming 锻炼了"在很大压力下把代码写出来写对"的能力——这个技能"到现在还是让我很受益"。

\subsection{竞赛编程对创业的影响}

虽然胡渊鸣没有显式提到竞赛编程与创业的关系，但从访谈中可以看到几个隐性连接：

\begin{enumerate}
\item \textbf{限时解题 $\rightarrow$ 快速迭代}：竞赛要求在有限时间内写出正确代码，创业要求快速试错、快速迭代。Meshy 从零到上线只用了8小时。
\item \textbf{边界条件思维 $\rightarrow$ 系统设计}：竞赛选手习惯考虑各种 corner case，这种思维在设计编译器和分布式系统时极为重要。
\item \textbf{竞争压力 $\rightarrow$ 抗压能力}：在高压环境下保持冷静和高效，是竞赛和创业的共同要求。
\item \textbf{排名文化 $\rightarrow$ 结果导向}：竞赛有明确的排名反馈，培养了"尊重事实、不自欺"的习惯。
\end{enumerate}

\subsection{清华姚班：认识天外有天}

进入清华姚班后，胡渊鸣收获的最大财富是"认识了一帮非常厉害的同学"。有人平时不怎么学但考试100分，有人轻松拿到 ACM World Final 金牌。他分享了一个印象深刻的挫败经历：

\begin{quote}
"第一门课叫'计算机应用数学'。一般到清华你会先学微积分、线性代数、概率论，但那门课好像直接 assume 你会了这些，直接开始讲。最离谱的是老师说'这个知识点我上课没来得及讲，所以我就把它出成了考试题目，希望你们在考试的时候能掌握这个知识点。'考完大家都说好难，有一个同学说'我要挂科了好难好难'——最后发现他考了98.5分。"
\end{quote}

\begin{importantbox}{更早认识厉害的人}
"更早地知道一些厉害的人，对一个人的成长有很大的帮助。它会让你知道天外有天、人外有人，原来自己也没有自己想的那么厉害，原来还有很多可以提升的地方。"——这种认知不是从姚班才开始的，早在初中搞竞赛、参加江苏省省队时就已经体验过。
\end{importantbox}

在清华时，中国做图形学的机会不多，最优秀的同学大多去搞 AI（当时正值 AlexNet、ResNet、AlphaGo 接连爆发）。胡渊鸣基本上自己在宿舍看 SIGGRAPH 论文、复现效果——这段独立研究的经历意外地训练了他后来创业时"从零开始"的能力。

他也去微软亚洲研究院（MSRA）实习，在 Steve Lin 老师的网络图形组做了 CVPR 论文（用神经网络预测白平衡）和 GAN+RL 的图像后处理工作——后者训练了1000多个模型才勉强得到一个可用结果。这段经历让他对当时 AI 技术的成熟度有了清醒认识：RL 的 sparse reward 问题"到今天其实都没有解决"——"今天也没有，我们过去一周捏了一周的 reward function。"

一个值得注意的细节是：他在清华本科期间所有能发表的论文\textbf{全部都是 AI 方向}的——"Graphics 一篇都没发出来"。虽然他最热爱 Graphics，但 AI 方向更容易出成果。这也预示了后来的趋势——Graphics 和 AI 的深度融合。

\begin{warningbox}{不要太迷信论文}
"科研有很多局限性，不能太迷信论文里面说的事情。"——他在复现 SIGGRAPH 论文时发现，很多论文并没有把所有细节写清楚，这个教训让他后来在太极项目中坚持"所有实验结果都可通过一行命令端到端复现"。
\end{warningbox}

\subsection{MIT 博士：从差点退学到创造太极}

2017年进入 MIT，胡渊鸣的第一年非常痛苦——与导师研究兴趣不一致，需要花大量时间做自己不 enjoy 的事情。他是一个"让我做 enjoy 的事可以1000\%努力，但做不 enjoy 的事会很痛苦"的人。

转机来自两位清华学长的建议：家俊建议他"把 simulator 搞成 differentiable"，这后来催生了 ChainQueen 和 DiffTaichi 两篇重要论文；俊彦则强调了"做有 impact 的 research"的重要性。最终他换到 Fredo Durand 和 Bill Freeman 的组，两位导师都比较 hands-off，给了他充分的研究自由。

\begin{importantbox}{如果重读 PhD 会做的四个改变}
\begin{enumerate}
\item \textbf{选择方向}：不去 NVIDIA 写太极，而是接受 Kaiming He 的 offer 去 FAIR 做 AI 研究——更早接触 AI frontier
\item \textbf{拓宽范围}：还是做太极，但不局限于 physical simulation，更早向 Transformer 优化方向发展
\item \textbf{高目标}：仍以一年毕业为目标——"求上而得中，求中得下，求下而不得"
\item \textbf{多交流}：花更多时间和不同领域的厉害的人交流，而不是天天把自己关在宿舍写太极
\end{enumerate}
\end{importantbox}

关于第四点，他引用了 Richard Hamming 的观点："如果你把自己关在办公室、把门关着每天做自己的事情，短时间内你当然能产出更多。但五年以后，你一定不知道该做什么。"——把门开着、多听其他人在做什么，才能发现值得解决的重要问题。

\subsection{MIT 与 Stanford：地理位置如何影响学术方向}

胡渊鸣对 MIT 和 Stanford 在 AI 时代的差异有一个独到观察。2017年他入学时，MIT 的 CSAIL 并没有很多人做 AI——整个 Boston 地区的创业氛围以 Biotech 为主，缺乏 AI 产业驱动。

\begin{knowledgebox}{产业生态如何塑造学术方向}
"MIT 和 Stanford 这种靠近产业中心的学校比起来，对 AI 的接受程度其实要低很多。由于没有产业驱动，整个 Boston 做的更多的还是 Bio。出厂公司基本上是 Bio 领域的。"——这导致 MIT 在整个 AI 时代跟得没有特别紧。"这可以是好事也可以是坏事——需要辩证地看。"
\end{knowledgebox}

这个观察对理解学术生态有重要启示：一所大学的研究方向不仅由教授们的兴趣决定，更受周围产业生态的深刻影响。Stanford 之所以成为 AI 重镇，与硅谷的科技产业密不可分。

\subsection{做科研的"95/5法则"}

胡渊鸣在知乎文章中引用过爱因斯坦的一句话："如果给我一个小时解决一个问题，我会花55分钟想什么是正确的问题，然后花5分钟去解决它。"这个"95/5法则"成为他对科研方法论的核心总结。

他在面试每一位 Meshy 求职者时都会问："你做的事情重要吗？"当被反问"你现在最近做的最重要的一件事情是什么"时，他的回答是：\textbf{招聘}（40\%的时间）、\textbf{思考未来战略}（30\%）、\textbf{参与日常业务}（30\%）。

\begin{warningbox}{大多数 PhD 在阶段二仍在"灌水"}
"你发了一篇论文，觉得做得很好，但和别人聊起来他们一点也不 exciting；你的文章只有几个人关注。为什么？因为你没有想清楚什么样的方向才是重要的方向。当你经过足够多的挫败和反思，才能知道：一个好的方向应该是什么样的。"——很多 PhD 在有了3-4篇顶会论文后，仍然在用阶段一的"卷产出"思维去灌水，而没有切换到阶段二的"寻找重要问题"模式。
\end{warningbox}

他以凯明（Kaiming He，何恺明）为例作为成功的"阶段二 PhD"典范——ResNet 就是"一个加号"，但它的影响力远超几十篇普通论文。他在知乎文章中写过"不要只做 SIGGRAPH 里面的小鱼，要去整个 CS 的池塘里面做大鱼"。如果太极 2.0 能在大模型时代有广泛应用、影响力达到1.0的50倍以上，他觉得"就做成一条大鱼了"。

\subsection{关于运气：99\%的归因}

胡渊鸣反复强调运气的重要性。当被问到"你很早就知道自己 enjoy 什么，这是天资还是运气"时，他的回答令人深思：

\begin{importantbox}{运气是底层操作系统}
"我觉得天资就是运气。你就是 genetically 抽中了一个比较好的基因，恰好在这个时代又需要你这个基因能做的事情。所以我觉得可能99\%都是运气。甚至你能努力也是运气的一部分——因为你有这个运气，所以你知道你喜欢什么，所以你可以努力。"
\end{importantbox}

这种对运气的高度归因并非消极——它实际上是 humble 的极致体现：承认自身优势的偶然性，对那些不如自己幸运的人保持同理心。同时这也解释了为什么他不批判选择"躺平"的人——"那是个人选择，能过得开心也很好。"

\subsection{本章小结}

胡渊鸣的成长路径展现了几个关键特质：（1）极早发现自己的兴趣并持续深耕；（2）通过实践"自己发现"物理定律，而非被动学习；（3）在每个阶段都遇到更强的人，保持 humble；（4）对学术研究保持批判性思维。他对导师关系的处理（从第一年的痛苦到换导师后的自由）、对职业选择的反思（后悔没去 FAIR）都展现了真诚的自我审视能力。

